% =============================================================================
% SECTION 1: INTRODUCTION
% =============================================================================
\section{Introduction}
\label{sec:introduction}

The emergence of Large Language Models (LLMs) as core components of software engineering
toolchains represents one of the most significant shifts in how software is designed, written,
maintained, and evolved. Systems such as GitHub Copilot, Amazon CodeWhisperer, Cursor, and a
growing ecosystem of research prototypes now assist developers with tasks ranging from
single-line autocomplete to full-feature implementation, test generation, bug localization, and
automated issue resolution. Yet beneath these impressive capabilities lies a deceptively hard
sub-problem: \emph{how should an LLM-based agent represent an entire software repository?}

A modern software repository is far more than a flat collection of source files. It encompasses
tens of thousands of functions distributed across hundreds of modules, intricate dependency
hierarchies, documented interfaces, test suites, configuration files, commit histories, issue
trackers, and implicit architectural conventions accumulated over years of collaborative
development. The informational content of a non-trivial open-source project routinely exceeds
ten million tokens---orders of magnitude beyond the context windows offered by even the most
capable contemporary LLMs. Representing this richness faithfully, efficiently, and in a form
that an LLM can exploit is the foundational challenge that all repository-aware AI systems must
address.

\subsection{Motivation and Problem Statement}

The importance of repository representation becomes immediately apparent when one examines the
failure modes of naive approaches. Early experiments with LLM-based code completion revealed
that models trained solely on individual files or functions failed catastrophically on
cross-file dependencies: they could not locate imported symbols, could not infer the contract
of a called function defined elsewhere in the repository, and could not respect project-specific
naming conventions that differ from open-source norms~\cite{ding2023crosscodeeval,
liu2024repobench}. The fundamental issue is one of context: the information needed to complete
a code fragment correctly is often not in the immediate vicinity of the cursor, but distributed
across the repository.

This observation has motivated a rapidly expanding body of research on repository-level software
engineering with LLMs. Researchers have proposed techniques ranging from static retrieval of
relevant snippets~\cite{zhang2023repocoder, lu2022reacc} to rich graph-based representations
of code structure~\cite{guo2021graphcodebert, liu2024graphcoder}, from hierarchical
summarization~\cite{dhulshette2025hierarchical, makharev2025code} to multi-step agentic
frameworks that dynamically navigate repository structure~\cite{yang2024sweagent,
zhang2024autocoderover, jimenez2024swebench}. This body of work has grown to encompass hundreds
of primary studies published in premier venues including ICSE, FSE, ASE, TOSEM, ACL, EMNLP,
ICLR, NeurIPS, and ICML.

Despite this explosion of work, the field lacks a principled taxonomy of \emph{repository
representation paradigms}. Existing surveys of LLM-based software engineering~\cite{hou2024llmse,
zhang2026survey, jiang2026survey} provide broad coverage of tasks and models but do not
systematically examine the representation dimension. Reviews of code
completion~\cite{husein2025codecompletion} focus on a single task and a single paradigm.
Reviews of retrieval-augmented generation for code~\cite{tao2026ragsurvey} treat retrieval as
the \emph{only} representation strategy, overlooking graph-based, embedding-based, agentic, and
summarization-based alternatives. This fragmented picture makes it difficult for both
researchers designing new systems and practitioners choosing deployment strategies to reason
about the design space in a principled way.

\subsection{The Repository Representation Problem}

We define \emph{repository representation} as the process by which an LLM-based agent
constructs and maintains a structured or unstructured encoding of the information contained in
a software repository, and makes that encoding available---directly or indirectly---to the
LLM's inference process. This definition is deliberately broad because the literature has
produced a diverse menagerie of strategies, each making different trade-offs among context
coverage, computational cost, inference latency, scalability, and task-specific performance.

The core challenge can be decomposed into three sub-problems. First, there is the
\emph{indexing problem}: given a repository of $N$ files containing $T$ tokens in total, how
does one construct a data structure that admits efficient queries of the form ``given the
current editing context $c$, find the $k$ most relevant fragments''? The answer depends
critically on what ``relevant'' means, which is itself task-dependent. For code completion, the
most relevant fragments are those most likely to contain symbols, types, or patterns needed
at the cursor. For bug localization, relevance is defined by causal proximity to the reported
failure. For code translation, relevance includes type definitions, library bindings, and
idiomatic patterns in the target language.

Second, there is the \emph{selection problem}: given an index of relevant fragments, how does
one select and rank a subset that fits within the LLM's context window? The context window is
simultaneously the most precious and most constrained resource in LLM-based software
engineering. Overfilling it with irrelevant context degrades generation quality and increases
inference cost; underfilling it wastes representational capacity and omits crucial information.
The selection problem is compounded by the well-documented phenomenon of ``lost in the middle'':
LLMs disproportionately attend to tokens at the beginning and end of their context window,
making the ordering of retrieved fragments non-trivially important~\cite{yusuf2025beyond}.

Third, there is the \emph{grounding problem}: once a representation has been selected and
placed in the LLM's context window, how does the model accurately interpret it? Different
representations---raw source code, API documentation, call graph adjacency lists, compressed
embeddings, natural language summaries---make different demands on the model's ability to
extract and apply the encoded information. A representation that is maximally information-dense
from an information-theoretic perspective may be opaque to a language model trained primarily
on natural language and syntactically structured source code.

\subsection{Research Questions}

This systematic literature review is organized around five research questions that together
cover the full scope of the repository representation problem:

\begin{description}[leftmargin=2.5em, labelwidth=2.5em]

\item[\textbf{\rqone}] \textit{What paradigms have been proposed for representing code
  repositories in LLM-based software engineering systems?} This question drives the
  construction of our seven-pillar taxonomy. We seek to enumerate and characterize the distinct
  representational strategies that appear in the literature, identify their key design choices,
  and trace their historical development from early dense embedding approaches through
  contemporary multi-agent frameworks.

\item[\textbf{\rqtwo}] \textit{How do repository representation paradigms map onto software
  engineering tasks?} The appropriateness of a given paradigm is not universal; different tasks
  impose different requirements on the representation. We investigate which paradigms are most
  commonly applied to which tasks, and whether there are systematic associations between paradigm
  choice and task characteristics such as locality, complexity, and interaction with the
  execution environment.

\item[\textbf{\rqthree}] \textit{What are the empirical performance characteristics of each
  paradigm, and how do they compare across standardized benchmarks?} We synthesize quantitative
  results from primary studies to characterize the performance envelope of each paradigm,
  identify which benchmarks are most widely used and thus best-suited for cross-study
  comparisons, and assess the degree to which benchmark results are reproducible and
  generalizable.

\item[\textbf{\rqfour}] \textit{What are the key trade-offs among repository representation
  paradigms along dimensions of context coverage, retrieval latency, training cost,
  scalability, and benchmark availability?} We go beyond raw performance to examine the
  operational characteristics that determine the suitability of each paradigm for real-world
  deployment, highlighting scenarios where each paradigm excels and where it fails.

\item[\textbf{\rqfive}] \textit{What are the open research challenges and future directions
  in repository representation for LLM-based software engineering?} We identify gaps in the
  current literature, propose concrete research agendas, and discuss emerging trends such as
  multi-paradigm hybridization, energy-aware evaluation, and adaptation to continuously evolving
  repositories.

\end{description}

\subsection{Contributions}

This paper makes the following primary contributions to the literature on LLM-based software
engineering:

\begin{enumerate}[leftmargin=*, label=(\arabic*)]

\item \textbf{Seven-Pillar PAR Taxonomy.} We present the first systematic taxonomy of
  repository representation paradigms (\parcmd{1}--\parcmd{7}), derived inductively from 160
  primary studies using thematic synthesis~\cite{cruzes2011thematic}. Each paradigm is
  characterized by its representational substrate, indexing mechanism, query interface,
  context management strategy, and typical deployment pattern.

\item \textbf{Bidirectional Task--Paradigm Mapping.} We construct a comprehensive mapping from
  the seven SE tasks most commonly addressed in the literature to the PAR taxonomy, and
  conversely from each PAR to the tasks it most effectively supports. This mapping exposes
  systematic under-coverage (tasks with few paradigm options) and over-coverage (paradigms
  applied uniformly to all tasks without empirical justification).

\item \textbf{Multi-Dimensional Trade-Off Analysis.} We evaluate each PAR across five
  operational dimensions---context coverage, retrieval latency, training cost, scalability,
  and benchmark availability---providing a structured basis for paradigm selection.

\item \textbf{Research Agenda.} We distill our findings into a prioritized set of seven open
  research challenges, each supported by evidence from the reviewed literature, and propose
  concrete research directions for each challenge.

\item \textbf{Replication Package.} We provide a publicly available replication package
  containing the full search protocol, screening decisions, data extraction spreadsheet, and
  synthesis code.

\end{enumerate}

\subsection{Scope and Boundaries}

This review focuses specifically on the \emph{representation} of code repositories, not on
LLM-based software engineering in general. We include studies that (a) target one or more
repository-level software engineering tasks, (b) use an LLM or a pre-trained code model as
the core inference engine, and (c) make explicit design choices about how repository content
is represented to the model. We exclude studies that focus exclusively on single-function or
single-file tasks with no cross-file context, studies that use LLMs purely as classifiers or
reviewers without generating or modifying code, and studies that focus on code understanding
for purposes other than software engineering (e.g., code plagiarism detection, malware analysis).

The temporal scope of this review is 2020--2026, capturing the period from the emergence of
transformer-based code models (CodeBERT~\cite{feng2020codebert}) through the latest
preprint literature available at the time of our search. This span encompasses both the
``pre-LLM'' era of specialized code representation models and the contemporary era of general
large language models such as GPT-4, Claude, and Llama.

\subsection{Paper Structure}

The remainder of this paper is organized as follows. Section~\ref{sec:methodology} describes
the systematic review methodology, including the search protocol, inclusion/exclusion criteria,
quality assessment, and data extraction procedures. Section~\ref{sec:taxonomy} presents the
seven-pillar PAR taxonomy, with detailed descriptions of each paradigm and representative
works. Section~\ref{sec:semapping} presents the bidirectional task--paradigm mapping and
analyzes coverage patterns. Section~\ref{sec:findings} answers our five research questions
with synthesized evidence from the primary studies. Section~\ref{sec:relatedwork} compares
this review with related surveys. Section~\ref{sec:threats} discusses threats to validity.
Section~\ref{sec:conclusion} concludes with a summary of key findings and the prioritized
research agenda. Appendices provide the full data extraction tables, quality assessment
results, and the complete list of primary studies.
